\section{Residual inertia and quantitative shape evolution}
\label{sec:r9-relative}

Let $N\ge1$, $m>1$, $s=m/2$, and let $a$ be a real, even,
positive elliptic homogeneous symbol of degree $m$, of class
$C^{N+2}$ away from the origin. Thus
$c|\xi|^m\le a(\xi)\le C|\xi|^m$.
Set $\alpha=2m/N$ and $p=2+2m/N$, and write
\[
 M(f)=\|f\|_2^2,\qquad
 A(f)=\int a(\xi)|\widehat f(\xi)|^2\,d\xi,\qquad
 E(f)=\tfrac12A(f)-\tfrac1p\|f\|_p^p,
 \qquad P(f)=\int\xi|\widehat f(\xi)|^2\,d\xi.
\]
We use the equivalent energy norm and its dual norm
\[
 \|f\|_X^2=M(f)+A(f),\qquad X=H^{m/2},\qquad
 \|h\|_{X^*}^2=\int\frac{|\widehat h(\xi)|^2}{1+a(\xi)}\,d\xi.
\]
Since $m>1$, differentiation maps $X$ continuously into $X^*$.
The map $f\mapsto |f|^\alpha f$ is locally Lipschitz from $X$ to
$X^*$, by H\"older's inequality and the embedding $X\hookrightarrow L^p$.
Pairings of real functions below are real dual pairings.

Let $M_c$ be the sharp critical mass, and let $\mathcal Q$ be the
set of real sharp optimizers normalized by $M(Q)=M_c$ and $A(Q)=1$.
Every $Q\in\mathcal Q$ satisfies
\[
 a(D)Q+\omega_*Q-|Q|^\alpha Q=0,
 \qquad \omega_*=\frac{m}{NM_c}.
\]
We use the optimizer compactness and phase coercivity established in
Lemmas~\ref{foundation:compact} and~\ref{foundation:phase}.
Their uniform consequences are recorded explicitly first.

\subsection{Uniform control near the optimizer set}

\begin{lemma}[Uniform optimizer bounds]\label{lem:r9rel-uniform}
There are $c_0,\kappa>0$, depending only on the fixed model, such that
for every $Q\in\mathcal Q$, the form
\[
 L_Q=a(D)+\omega_*-|Q|^\alpha
\]
obeys
\begin{align}
 \langle b,L_Qb\rangle&\ge c_0\|b\|_X^2
 &&\text{for real }b\in X,\ \langle b,Q\rangle=0,
 \label{eq:r9rel-uniformcoercivity}\\
 \|\zeta\cdot\nabla Q\|_{X^*}&\ge\kappa|\zeta|
 &&\text{for }\zeta\in\mathbb R^N,
 \label{eq:r9rel-derivativeinjective}\\
 \|\tau Q+i\zeta\cdot\nabla Q\|_{X^*}
 &\ge\kappa(|\tau|+|\zeta|)
 &&\text{for }(\tau,\zeta)\in\mathbb R\times\mathbb R^N.
 \label{eq:r9rel-parameterinjective}
\end{align}
\end{lemma}
\begin{proof}
The normalized optimizer set is compact in $X$ modulo translations.
Moreover, $Q\mapsto |Q|^\alpha$ is continuous in the form norm
$X\to X^*$. If \eqref{eq:r9rel-uniformcoercivity} failed, translate
a sequence of optimizers to obtain $Q_n\to Q$ strongly in $X$, and
choose real $b_n\perp Q_n$ with $\|b_n\|_X=1$ and
$\langle b_n,L_{Q_n}b_n\rangle\to0$.
The projection of $b_n$ onto $Q$ tends to zero. Form convergence and
the strict coercivity for this fixed $Q$ contradict the unit norm.

For a nonzero direction $\zeta$, the identity
$\zeta\cdot\nabla Q=0$ would give
$(\zeta\cdot\xi)\widehat Q(\xi)=0$.
An $L^2$ function supported on a hyperplane is zero, contrary to
$M(Q)=M_c$. Compactness modulo translations and continuity of
differentiation from $X$ to $X^*$ make this lower bound uniform over
$Q$ and over the unit sphere of directions.
Finally, real and purely imaginary functions are orthogonal for
the real-even multiplier defining the $X^*$ inner product.
The norms $\|Q\|_{X^*}$ have a uniform positive lower bound by
optimizer compactness, which proves
\eqref{eq:r9rel-parameterinjective}.
\end{proof}

\begin{lemma}[Uniform proximity and phase choice]\label{lem:r9rel-proximity}
There is a modulus $\varepsilon(\delta)\to0$ as $\delta\downarrow0$
such that
\[
 A(U)=1,\qquad E(U)<0,\qquad M_c<M(U)<M_c+\delta
\]
imply that, after a translation and a constant phase, some
$Q\in\mathcal Q$ satisfies
\[
 \|U-Q\|_X\le\varepsilon(\delta),\qquad
 \langle\operatorname{Im}U,Q\rangle=0.
\]
\end{lemma}
\begin{proof}
Failure of the first assertion would contradict normalized
near-optimizer compactness. A further phase tending to zero makes
$\int UQ$ real and positive; this is possible because the integral
tends to $M_c$. It changes the proximity modulus only by a fixed
factor and gives the asserted orthogonality.
\end{proof}

\subsection{A quantitative residual theorem}

\begin{theorem}[Positive inertia with a residual error]
\label{thm:r9rel-residual}
There exist $\delta,c,C>0$, depending only on $N,m,a$, such that
the following holds. Let $U\in X$ satisfy
\[
 M_c<M(U)<M_c+\delta,\qquad E(U)<0.
\]
For arbitrary $\omega\in\mathbb R$ and $v\in\mathbb R^N$, define
\[
 F=a(D)U+\omega U+iv\cdot\nabla U-|U|^\alpha U\in X^*.
\]
Then
\begin{equation}\label{eq:r9rel-residualinertia}
 v\cdot P(U)\ge cA(U)^{2/m-1}|v|^2
 -C\int\frac{|\widehat F(\xi)|^2}{A(U)+a(\xi)}\,d\xi.
\end{equation}
The constants are independent of frequency and velocity.
No radial symmetry or spatial moment is required.

If, in addition, the scalar dilation defect satisfies
\begin{equation}\label{eq:r9rel-defecthyp}
 D:=mE(U)-\tfrac12v\cdot P(U)
 \ge-\theta m|E(U)|,
 \qquad 0\le\theta<1,
\end{equation}
then, for some $c_\theta>0$ depending only on the model and $\theta$,
\begin{equation}\label{eq:r9rel-residuallower}
 \int\frac{|\widehat F(\xi)|^2}{A(U)+a(\xi)}\,d\xi
 \ge c_\theta\bigl(|E(U)|+A(U)^{2/m-1}|v|^2\bigr).
\end{equation}
\end{theorem}
\begin{proof}
First normalize $A(U)=1$, and put $\rho=\|F\|_{X^*}$.
Use Lemma~\ref{lem:r9rel-proximity} to arrange proximity to a real
$Q\in\mathcal Q$ and the stated phase orthogonality. All quantities
in the conclusion are invariant under these translations and phases.

Write $G(U)=a(D)U-|U|^\alpha U$. Its $X^*$ norm is bounded uniformly
in the normalized mass window. Replacing $Q$ by $U$ in
\eqref{eq:r9rel-parameterinjective} costs at most
$C\varepsilon(\delta)(|\omega|+|v|)$.
For sufficiently small $\delta$, absorption in
$F=G(U)+\omega U+iv\cdot\nabla U$ therefore gives
\begin{equation}\label{eq:r9rel-parameterbound}
 |\omega|+|v|\le C(1+\rho).
\end{equation}
Subtracting $G(Q)=-\omega_*Q$, and using the local Lipschitz
continuity of $G$, now gives
\begin{equation}\label{eq:r9rel-parametersmall}
 |\omega-\omega_*|+|v|
 \le C\{\rho+\varepsilon(\delta)(1+\rho)\}.
\end{equation}

Fix a sufficiently small $\rho_0>0$, and then shrink $\delta$.
For $\rho\le\rho_0$, the form
$L_U=a(D)+\omega-|U|^\alpha$ is uniformly close to $L_Q$.
Consequently,
\begin{equation}\label{eq:r9rel-perturbedcoercivity}
 \langle b,L_Ub\rangle\ge c_1\|b\|_X^2
 \quad\text{for real }b\perp Q,
\end{equation}
with $c_1>0$ independent of $U,\omega,v$ in this regime.
Write $U=h+ib$. Its imaginary residual equation is
\[
 \operatorname{Im}F=L_Ub+v\cdot\nabla h.
\]
The identity $P(U)=-2\langle b,\nabla h\rangle$ is valid by $m>1$.
Thus
\begin{align}
 \tfrac12v\cdot P(U)
 &=\langle b,L_Ub\rangle-\langle\operatorname{Im}F,b\rangle
 \notag\\
 &\ge\tfrac12c_1\|b\|_X^2-C\rho^2.
 \label{eq:r9rel-phasepairing}
\end{align}
The imaginary equation also gives, by
\eqref{eq:r9rel-derivativeinjective},
\[
 \kappa|v|\le\|v\cdot\nabla Q\|_{X^*}
 \le\rho+C\|b\|_X+C\varepsilon(\delta)|v|.
\]
Here $L_U:X\to X^*$ has uniformly bounded norm.
Absorbing the last term yields
$|v|\le C(\rho+\|b\|_X)$. Combining this with
\eqref{eq:r9rel-phasepairing} proves
\begin{equation}\label{eq:r9rel-normalized}
 v\cdot P(U)\ge c|v|^2-C\rho^2
\end{equation}
when $\rho\le\rho_0$.

If $\rho>\rho_0$, normalized mass and kinetic energy give
$|P(U)|\le C$, while \eqref{eq:r9rel-parameterbound} gives
$|v|\le C(1+\rho)$. Hence
\[
 c|v|^2-v\cdot P(U)\le C(1+\rho)^2
 \le C_{\rho_0}\rho^2.
\]
After enlarging $C$, \eqref{eq:r9rel-normalized} holds for every
residual size.

For general $A=A(U)>0$, set $\lambda=A^{-1/m}$ and
\[
 U_\lambda(x)=\lambda^{N/2}U(\lambda x),\qquad
 \omega_\lambda=\lambda^m\omega,\qquad
 v_\lambda=\lambda^{m-1}v.
\]
Then $A(U_\lambda)=1$,
$v_\lambda\cdot P(U_\lambda)=A^{-1}v\cdot P(U)$, and the
residual scales as
$F_\lambda(x)=\lambda^{N/2+m}F(\lambda x)$.
A change of Fourier variables gives exactly
\[
 \|F_\lambda\|_{X^*}^2
 =A^{-1}\int\frac{|\widehat F(\xi)|^2}{A+a(\xi)}\,d\xi.
\]
Rescaling \eqref{eq:r9rel-normalized} proves
\eqref{eq:r9rel-residualinertia}. Finally,
\eqref{eq:r9rel-defecthyp} implies
\[
 v\cdot P(U)=2mE(U)-2D
 \le-2(1-\theta)m|E(U)|.
\]
Substitution into \eqref{eq:r9rel-residualinertia} proves
\eqref{eq:r9rel-residuallower}.
\end{proof}

\begin{remark}
The defect $D$ is defined without a spatially weighted pairing of
$F$. When that pairing is justified,
$D=\operatorname{Re}\langle F,(x\cdot\nabla+N/2)U\rangle$;
the same conclusion holds through the corresponding cutoff limit.
For an exact traveling profile the Pohozaev identity gives $D=0$.
The residual theorem therefore refines the profile exclusion by
quantifying how far an approximate relative equilibrium with small
dilation defect must remain from satisfying its profile equation.
\end{remark}

\subsection{The positive response tensor}

Let $L_Q^{-1}$ denote the coercive form inverse on the real
$L^2$-orthogonal complement of $Q$. Since
$\langle\partial_jQ,Q\rangle=0$, the matrix
\begin{equation}\label{eq:r9rel-tensor}
 T_{Q,ij}=\langle\partial_iQ,L_Q^{-1}\partial_jQ\rangle
\end{equation}
is well defined.

\begin{proposition}[Uniform effective inertia]\label{prop:r9rel-tensor}
The matrices $T_Q$ are symmetric positive definite, with eigenvalues
bounded above and below by positive model-dependent constants,
uniformly for $Q\in\mathcal Q$.
Suppose
\[
 A(U_n)=1,\quad E(U_n)<0,\quad M(U_n)\downarrow M_c,
 \qquad \rho_n=\|F_n\|_{X^*}\longrightarrow0,
\]
where $F_n$ uses arbitrary parameters $(\omega_n,v_n)$.
For the optimizers $Q_n$ supplied by
Lemma~\ref{lem:r9rel-proximity},
\begin{equation}\label{eq:r9rel-tensorresponse}
 P(U_n)=2T_{Q_n}v_n+o(|v_n|)+O(\rho_n),
\end{equation}
with a uniform constant in the $O$ term.
\end{proposition}
\begin{proof}
The $L^2$ projection onto $Q^\perp$ is uniformly bounded on $X$
and its dual. Thus the restricted dual norm is uniformly equivalent
to the full norm for functionals annihilating $Q$.
For $d=\zeta\cdot\nabla Q$, uniform coercivity and boundedness of
$L_Q$ give
\[
 c\|d\|_{X^*}^2\le\langle d,L_Q^{-1}d\rangle
 \le C\|d\|_{X^*}^2.
\]
Lemma~\ref{lem:r9rel-uniform} proves both eigenvalue bounds.

Arrange $U_n=h_n+ib_n$ as in the residual proof.
From \eqref{eq:r9rel-perturbedcoercivity} and the exact momentum
pairing one also obtains
\[
 c_1\|b_n\|_X^2
 \le\tfrac12v_n\cdot P(U_n)+\rho_n\|b_n\|_X
 \le C|v_n|\|b_n\|_X+\rho_n\|b_n\|_X.
\]
Hence $\|b_n\|_X\le C(|v_n|+\rho_n)$.
Using the imaginary residual equation, form convergence, and
\eqref{eq:r9rel-parametersmall},
\[
 L_{Q_n}b_n=-v_n\cdot\nabla Q_n
 +o_{X^*}(|v_n|)+O_{X^*}(\rho_n).
\]
Project onto $Q_n^\perp$ and invert $L_{Q_n}$.
Pairing the resulting identity with $-2\nabla h_n$, and using
$\|h_n-Q_n\|_X\to0$, proves
\eqref{eq:r9rel-tensorresponse}.
\end{proof}

In particular, a residual satisfying $\rho_n=o(|v_n|)$ forces
strictly positive velocity--momentum pairing at leading order.
This conclusion uses neither uniqueness of the optimizer nor
nondegeneracy of the real-part linearized operator.

\subsection{Relative-periodic solutions and shape speed}

The following statements impose explicit regularity sufficient for
the differentiated cutoff virial identity. They do not assume that
every energy-space solution has this additional regularity.

\begin{theorem}[Nonvanishing shape speed]\label{thm:r9rel-periodic}
Suppose a solution of
\[
 iu_t=a(D)u-|u|^\alpha u
\]
has the relative-periodic representation
\[
 u(t,x)=e^{i\omega t}V(t,x-vt),\qquad V(t+T)=V(t),
\]
where $T>0$ and
\begin{equation}\label{eq:r9rel-periodicregularity}
 V\in C_{\mathrm{per}}([0,T];H^{m+1})
 \cap C^1_{\mathrm{per}}([0,T];H^{m/2}).
\end{equation}
Assume that the conserved energy is negative and that
$M_c<M<M_c+\delta$, with $\delta$ as in
Theorem~\ref{thm:r9rel-residual}. Then
\begin{equation}\label{eq:r9rel-periodicpohozaev}
 v\cdot P=2mE<0.
\end{equation}
At every time and for every real number $\nu$,
\begin{equation}\label{eq:r9rel-shapespeed}
 \int\frac{|\widehat{\partial_tV-i\nu V}(t,\xi)|^2}
 {A(V(t))+a(\xi)}\,d\xi
 \ge c\bigl(|E|+A(V(t))^{2/m-1}|v|^2\bigr),
\end{equation}
where $c>0$ depends only on the model.
\end{theorem}
\begin{proof}
The equation in this frame is
\[
 iV_t=a(D)V+\omega V+iv\cdot\nabla V-|V|^\alpha V.
\]
Choose a real $\chi\in C_c^\infty$ equal to one near the origin,
and set
\[
 f_R(x)=x\chi(x/R),\qquad
 L_R=f_R\cdot\nabla+\tfrac12\operatorname{div}f_R,
 \qquad J_R(V)=\operatorname{Im}\int\overline V L_RV.
\]
The operator $L_R$ is skew-symmetric. Differentiation gives
\begin{align}
 J_R'(t)
 &=\langle V,[a(D),L_R]V\rangle
 +\langle V,[iv\cdot\nabla,L_R]V\rangle\notag\\
 &\quad-2\left(\tfrac12-\tfrac1p\right)
 \int\operatorname{div}f_R\,|V|^p.
 \label{eq:r9rel-cutoffvirial}
\end{align}
The frequency term vanishes by skew symmetry.
Assumption~\eqref{eq:r9rel-periodicregularity} justifies the
differentiation and all pairings. In particular,
$|V|^\alpha V\in L^2$ by Sobolev embedding from $H^{m+1}$,
and $\nabla V\in L^2$, so the nonlinear integration by parts is valid.

For completeness, the kinetic commutator has Fourier kernel,
up to the fixed normalization factor,
\[
 \frac i2\bigl(a(\xi)-a(\eta)\bigr)
 (\xi+\eta)\cdot\widehat f_R(\xi-\eta).
\]
Homogeneity gives
$|a(\xi)-a(\eta)|\le C|\xi-\eta|(|\xi|+|\eta|)^{m-1}$.
After division by the two energy weights, the kernel is bounded
by $C|k|\langle k\rangle^{m/2}|\widehat f_R(k)|$,
where $k=\xi-\eta$. Since
$\widehat f_R(k)=R^{N+1}\widehat f(Rk)$, this majorant has
uniformly bounded $L^1$ norm for $R\ge1$.
Young's inequality gives a uniform bound $X\to X^*$.
For Schwartz functions, Fourier integration by parts and
$\xi\cdot\nabla a=ma$ give the limiting form $mA$.
Density and the uniform bound extend the limit to every $V\in X$.
The identical degree-one argument for $iv\cdot\nabla$ gives
the limit $-v\cdot P$. Finally, dominated convergence gives the
nonlinear limit $-2m\|V\|_p^p/p$.
Thus the right side of \eqref{eq:r9rel-cutoffvirial} tends to
\[
 mA(V)-v\cdot P-\frac{2m}{p}\|V\|_p^p
 =2mE-v\cdot P.
\]
Uniform form bounds and boundedness of the trajectory on the period
permit dominated convergence in time. Integrate first over
$[0,T]$. Periodicity gives $J_R(T)-J_R(0)=0$ for every $R$;
then let $R\to\infty$. This proves
\eqref{eq:r9rel-periodicpohozaev} without a spatial moment.

At a fixed time apply Theorem~\ref{thm:r9rel-residual} to $V(t)$
with frequency $\omega+\nu$ and velocity $v$.
Its residual is
\[
 F=iV_t+\nu V=i(V_t-i\nu V).
\]
Equation~\eqref{eq:r9rel-periodicpohozaev} gives $D=0$.
The lower bound \eqref{eq:r9rel-residuallower}, with $\theta=0$,
is precisely \eqref{eq:r9rel-shapespeed}.
\end{proof}

The result also applies to a profile periodic up to a constant
phase: a linear phase adjustment makes the representation periodic,
and the cutoff functional is phase invariant.

\begin{corollary}[Moving-center cost over a cycle]
\label{cor:r9rel-movingcenter}
Under the same mass and energy hypotheses, suppose instead that
\[
 W(t,x)=e^{i\gamma(t)}u(t,x+c(t)),\qquad W(t+T)=W(t),
\]
where $c,\gamma$ are $C^1$ and $W$ has the regularity
\eqref{eq:r9rel-periodicregularity}.
Put $d=c(T)-c(0)$. Then
\begin{equation}\label{eq:r9rel-netdisplacement}
 P\cdot d=2mET,
\end{equation}
and
\begin{align}
 &\int_0^T\int
 \frac{|\widehat{W_t}(t,\xi)|^2}{A(W(t))+a(\xi)}\,d\xi\,dt
 \notag\\[-2pt]
 &\hspace{15mm}\ge
 c\left\{|E|T+
 \int_0^T A(W(t))^{2/m-1}|c'(t)|^2\,dt\right\}.
 \label{eq:r9rel-cyclecost}
\end{align}
The second term in braces is at least
\begin{equation}\label{eq:r9rel-displacementcost}
 \frac{|d|^2}{\displaystyle\int_0^T A(W(t))^{1-2/m}\,dt}.
\end{equation}
\end{corollary}
\begin{proof}
The equation for $W$ is
\[
 iW_t=a(D)W-\gamma'(t)W+ic'(t)\cdot\nabla W-|W|^\alpha W.
\]
The same cutoff virial calculation, integrated over the period,
gives \eqref{eq:r9rel-netdisplacement}.
At each time, Theorem~\ref{thm:r9rel-residual} applies with
frequency $-\gamma'(t)$, velocity $c'(t)$, and residual $iW_t$.
Integrating \eqref{eq:r9rel-residualinertia}, and substituting
\eqref{eq:r9rel-netdisplacement}, proves
\eqref{eq:r9rel-cyclecost}. Weighted Cauchy--Schwarz gives
\eqref{eq:r9rel-displacementcost}.
\end{proof}

These inequalities give necessary quantitative restrictions on
near-threshold negative-energy coherent dynamics. They exclude
vanishing shape speed in the stated relative-periodic class, even
after a time-dependent choice of center and phase. They neither
exclude all breathers nor establish a mass gap for arbitrary bounded
global solutions. The left side of \eqref{eq:r9rel-cyclecost} measures
evolution in a dual energy norm; it is not a physical dissipation,
and no loss of conserved energy is asserted.
